\documentclass[11pt]{article}

\usepackage[margin=1in]{geometry}
\usepackage{amsmath,amssymb,amsthm,mathtools}
\usepackage{booktabs,longtable,array}
\usepackage{microtype}
\usepackage[hidelinks]{hyperref}
\usepackage{xcolor}

\newtheorem{theorem}{Theorem}[section]
\newtheorem{lemma}[theorem]{Lemma}
\newtheorem{corollary}[theorem]{Corollary}
\newtheorem{proposition}[theorem]{Proposition}
\theoremstyle{definition}
\newtheorem{remark}[theorem]{Remark}

\newcommand{\vp}{v_p}
\newcommand{\Z}{\mathbb Z}
\newcommand{\Q}{\mathbb Q}
\newcommand{\CT}{\operatorname{CT}}
\newcommand{\OEIS}[1]{\href{https://oeis.org/#1}{#1}}

\title{Elementary Transfer Principles for Supercongruence Towers\\
\large A compact portfolio draft for external review}
\author{Ravi Andrew Bajaj \and Alexander Burns}
\date{August 1, 2026}

\begin{document}
\maketitle

\begin{center}
\fbox{\parbox{0.91\textwidth}{\small
Working draft.  The arguments are complete proof candidates with exact
transcription checks, but the manuscript has not been peer reviewed or
submitted.  The purpose of this version is focused mathematical review,
especially by Peter Bala, Paul D. Hanna, and specialists in $p$-adic
supercongruences.}}
\end{center}

\begin{abstract}
We collect four elementary mechanisms for adjacent prime-power
supercongruences.  First, a carry-and-transfer argument proves the
all-prime congruence conjectured for \OEIS{A183068}.  Second, complete
unit blocks prove a cubic transfer theorem for balanced products of gamma
values at rational slopes.  This theorem settles the congruence part of a
packet of fractional-factorial conjectures and, together with known or
finite-sum integrality, proves the full towers for \OEIS{A364175}, every
row of \OEIS{A365025}, and every homogeneous row of \OEIS{A364513}.
Third, the same unit-block method gives uniform half-integral binomial
families containing \OEIS{A091527}, \OEIS{A262732}, \OEIS{A275652}, and
\OEIS{A275654}.  Fourth, a complementary-factor decomposition of
MacMahon's product proves a nonlinear $p^{4r}$ Frobenius tower for plane
partitions in every $N\times N\times cN$ box, including \OEIS{A008793},
\OEIS{A352656}, and \OEIS{A352657}.  We separate theorem, integrality,
finite verification, and literature-priority status throughout.  A final
section packages the towers as arithmetic state sums and explains the
precise, limited analogy with spin-foam coarse graining.
\end{abstract}

\tableofcontents

\section{Introduction and status boundary}

An adjacent prime-power supercongruence compares the values of an integer
sequence at two successive points of a $p$-adic tower.  The typical form is
\[
 a(np^r)\equiv a(np^{r-1})\pmod {p^{\lambda r}},
\]
while product sequences can obey the nonlinear variant
\[
 a(np^r)\equiv a(np^{r-1})^p\pmod {p^{\lambda r}}.
\]
The main theme here is that many apparently unrelated OEIS conjectures are
controlled by one of two local operations:
\begin{enumerate}
\item discard summands whose base-$p$ carries already supply the modulus,
      and rescale the remaining summands; or
\item split a product into its $p$-divisible factors and complete blocks of
      $p$-adic units, then annihilate the first two reciprocal moments.
\end{enumerate}

The source conjectures were formulated by Peter Bala on records created by
several OEIS contributors.  In particular, Paul D. Hanna created
\OEIS{A183068}, and Bala supplied its factorial-sum formula and
supercongruence conjecture.  The present manuscript attributes the source
statements but does not assign authorship to their originators without
their consent.

There are three separate questions for every claimed application.
\begin{enumerate}
\item Is the displayed value an integer (or at least $p$-integral)?
\item Does the adjacent quotient have the asserted $p$-adic precision?
\item Is the result new in the literature?
\end{enumerate}
The proofs below address the first two only where stated.  Negative searches
are not treated as priority certificates.

\section{Unit blocks and classical binomial transfer}

We write $v_p(x)$ for the normalized $p$-adic valuation of a nonzero
rational number.  A congruence between rational numbers is interpreted in
$\Q_p$.

\begin{lemma}[complete unit block]\label{lem:unit}
Let $p\geq5$, let $P=p^e$ with $e\geq1$, and let $I$ be any union of
complete intervals of length $P$.  Summing only over integers prime to $p$,
\[
 \sum_{j\in I,\,p\nmid j}\frac1j\equiv0\pmod {P^2},
 \qquad
 \sum_{j\in I,\,p\nmid j}\frac1{j^2}\equiv0\pmod P.
\]
\end{lemma}

\begin{proof}
For a single reduced residue block, multiplication by a unit $c$ with
$c^2\not\equiv1\pmod p$ permutes the units.  Hence the reciprocal-square
sum is congruent to $c^{-2}$ times itself and is zero modulo $P$.
Pairing $u$ and $P-u$ gives
\[
 \frac1u+\frac1{P-u}\equiv-\frac P{u^2}\pmod {P^2},
\]
so the reciprocal sum vanishes modulo $P^2$.  Translation by $tP$ preserves
the two precisions because
\[
 (tP+u)^{-1}\equiv u^{-1}-tPu^{-2}\pmod {P^2},\qquad
 (tP+u)^{-2}\equiv u^{-2}\pmod P.
\]
Adding blocks proves the result.
\end{proof}

We also use the classical Ljunggren--Jacobsthal--Kazandzidis congruence in
the following adjacent form: if $p\geq5$, $N=pM$, and $e=v_p(N)$, then
\begin{equation}\label{eq:LJK}
 \frac{\binom{aN}{bN}}{\binom{aM}{bM}}
 \equiv1\pmod {p^{3e}}
\end{equation}
whenever the binomial coefficients are defined.  Refined versions and the
small-prime losses are discussed in \cite{Kazandzidis,StraubLjunggren,Zudilin}.

\section{Balanced rational gamma ratios}

\subsection{A rational-binomial transfer}

For positive integers $a,b,q$ with $a>bq$, define
\[
 G_{a,b,q}(N)=\binom{aN/q}{bN}
 =\frac{\Gamma(aN/q+1)}{\Gamma(bN+1)\Gamma((a-bq)N/q+1)}.
\]

\begin{theorem}[rational-binomial transfer]\label{thm:rbin}
Let $p\geq5$, $p\nmid q$, and $N=pM$.  If $e=v_p(N)$, then
\[
 \frac{G_{a,b,q}(N)}{G_{a,b,q}(M)}\equiv1\pmod {p^{3e}}.
\]
\end{theorem}

\begin{proof}
The finite product
\[
 G_{a,b,q}(N)=\frac1{q^{bN}(bN)!}
 \prod_{j=0}^{bN-1}(aN-qj)
\]
has a numerator factor divisible by $p$ exactly when $p\mid j$.  Cancelling
those factors and the corresponding denominator factors against the level
$M$ product gives the exact identity
\[
 \frac{G_{a,b,q}(N)}{G_{a,b,q}(M)}
 =\prod_{\substack{1\leq j<bN\\p\nmid j}}
 \left(1-\frac{aN}{qj}\right).
\]
The sign disappears because $b(N-M)=bM(p-1)$ is even.  Put $P=p^e$.
The index interval is a union of complete $P$-blocks.  Lemma
\ref{lem:unit} shows that the linear and quadratic elementary-symmetric
terms are divisible by $P^3$.  Every term of degree at least three contains
$N^3$, while all denominators are $p$-adic units.  Thus the product is
$1$ modulo $P^3$.
\end{proof}

\subsection{The residue-balanced theorem}

Let $E:\Q_{>0}\to\Z$ have finite support and put
\[
 F(N)=\prod_{\alpha\in\Q_{>0}}\Gamma(\alpha N+1)^{E(\alpha)}.
\]
Assume
\begin{equation}\label{eq:slope}
 \sum_\alpha \alpha E(\alpha)=0
\end{equation}
and, for every nonzero class $\rho\in\Q/\Z$,
\begin{equation}\label{eq:residue}
 \sum_{\alpha\equiv\rho\pmod\Z}E(\alpha)=0.
\end{equation}
Let $Q$ be a common denominator of the slopes.

\begin{theorem}[residue-balanced gamma transfer]\label{thm:gamma}
For $p\geq5$, $p\nmid Q$, and $n,r\geq1$,
\[
 \frac{F(np^r)}{F(np^{r-1})}\equiv1\pmod {p^{3r}}.
\]
If $F(np^{r-1})$ is $p$-integral, then
\[
 F(np^r)\equiv F(np^{r-1})\pmod {p^{3r}}.
\]
\end{theorem}

\begin{proof}
Condition \eqref{eq:residue} pairs every numerator gamma factor of
nonintegral slope with a denominator factor in the same residue class.
For $\alpha>\beta$ and $d=\alpha-\beta\in\Z_{>0}$,
\[
 \frac{\Gamma(\alpha N+1)}{\Gamma(\beta N+1)}
 =\binom{\alpha N}{dN}\Gamma(dN+1).
\]
Use the inverse identity if $\alpha<\beta$.  Repeating this operation leaves
rational binomials governed by Theorem \ref{thm:rbin} and a balanced
integer-slope factorial ratio.  A balanced integer-slope ratio factors as a
Laurent product of $\binom{jN}{N}$, so \eqref{eq:LJK} applies to every
factor.  Products and inverses preserve the congruence.  Multiplication by
the lower $p$-integral value proves the final assertion.
\end{proof}

\subsection{Applications}

The sequence \OEIS{A364175} is
\[
 A(N)=\frac{(6N)!(2N/3)!}{(3N)!(2N)!(5N/3)!}
 =\binom{6N}{3N}\binom{3N}{N}\binom{5N/3}{N}^{-1}.
\]
David Radcliffe proved its integrality \cite{Radcliffe}.  Theorem
\ref{thm:gamma} therefore gives the following complete tower.

\begin{corollary}\label{cor:a364175}
For every prime $p\geq5$ and $n,r\geq1$,
\[
 A(np^r)\equiv A(np^{r-1})\pmod {p^{3r}}.
\]
\end{corollary}

For a fixed $s\geq0$, define
\[
 U_s(N)=
 \frac{\Gamma(N/2+1)\Gamma((2s+1)N+1)\Gamma((2s+1/2)N+1)}
 {\Gamma(N+1)\Gamma(sN+1)^2\Gamma((s+1/2)N+1)^2}.
\]
The OEIS identity
\[
 U_s(N)=\sum_{j=0}^{sN}\binom{(2s+1)N}{sN-j}^2
                   \binom{N+j-1}{j}
\]
proves integrality.  The residue class $1/2+\Z$ has two numerator and two
denominator factors, so Theorem \ref{thm:gamma} proves the conjectured
cubic tower for every row of \OEIS{A365025}, including \OEIS{A365026} and
\OEIS{A365027}.

For $s\geq3$, the row family on \OEIS{A364513} has the gamma form
\[
 V_s(N)=\frac{s-2}{s}
 \frac{\Gamma((s+2)N+1)\Gamma(sN/2+1)^2}
 {\Gamma((s+2)N/2+1)\Gamma(sN+1)
  \Gamma((s-2)N/2+1)\Gamma(N+1)^2}
\]
and the integral finite sum
\[
 V_s(N)=\sum_{j=0}^N\binom{sN-1}{N-j}^2
                     \binom{(s-2)N+j-2}{j}.
\]
The constant $(s-2)/s$ cancels from adjacent quotients.  Theorem
\ref{thm:gamma} proves the tower for every fixed $s\geq3$, including
\OEIS{A364515}, \OEIS{A364516}, and \OEIS{A364517}.  Row $1$ is not
homogeneous and is not claimed here.

The same residue test applies to the displayed gamma quotients on
\OEIS{A364172}--\OEIS{A364184}.  It proves their cubic transfer in $\Q_p$;
an integer congruence still requires global integrality.  This logical
separation is essential.

\section{The all-prime A183068 theorem}

Define
\[
 a(N)=\sum_{k=0}^N F(N,k),\qquad
 F(N,k)=\frac{(2N+2k)!}{k!^4(N-k)!^2}.
\]
The summand is the six-part multinomial coefficient
\[
 F(N,k)=\binom{2N+2k}{k,k,k,k,N-k,N-k}.
\]

\begin{lemma}[carry depth]\label{lem:carry}
Suppose $p^t\mid N$ and $s=v_p(k)<t$.  Then
\[
 v_p(F(N,k))\geq2(t-s).
\]
For $p=2$ the lower bound improves to $2(t-s)+1$.
\end{lemma}

\begin{proof}
At level $q=p^i$, write $N=qM$ and $k=qa+u$ with $0<u<q$.  Legendre's
formula gives the level contribution
\[
 \left\lfloor\frac{2N+2k}{q}\right\rfloor
 -4\left\lfloor\frac{k}{q}\right\rfloor
 -2\left\lfloor\frac{N-k}{q}\right\rfloor
 =2+\left\lfloor\frac{2u}{q}\right\rfloor\geq2.
\]
There are $t-s$ active levels.  If $p=2$, then at $i=s+1$ one has
$u=2^s$, and this level contributes $3$.
\end{proof}

We use the multinomial form of the prime-power
Ljunggren--Jacobsthal--Kazandzidis estimate.  If every positive component
$b_i$ is divisible by $p^s$, then
\begin{equation}\label{eq:multiscale}
 \frac{\binom{p\sum b_i}{pb_1,\ldots,pb_m}}
      {\binom{\sum b_i}{b_1,\ldots,b_m}}
 \equiv1\pmod {p^{3(s+1)-\epsilon_p}},
\end{equation}
where $\epsilon_2=2$, $\epsilon_3=1$, and $\epsilon_p=0$ for $p\geq5$.
The possible binary sign is harmless here: for $s\geq1$ its exceptional
parity condition cannot occur, while at $s=0$ the modulus is $2$.

\begin{lemma}[summand transfer]\label{lem:summand}
Let $N'=np^{r-1}$ and $0\leq\ell\leq N'$.  Then
\[
 F(pN',p\ell)\equiv F(N',\ell)\pmod {p^{2r}}.
\]
\end{lemma}

\begin{proof}
Apply \eqref{eq:multiscale} to
\[
 \ell,\ell,\ell,\ell,N'-\ell,N'-\ell.
\]
Let $s$ be the smallest valuation among positive components.  If
$s<r-1$, Lemma \ref{lem:carry} contributes $2(r-1-s)$ powers (one more at
$p=2$), while \eqref{eq:multiscale} contributes
$3(s+1)-\epsilon_p$.  Their sum is at least $2r$.  If $s\geq r-1$,
the scaling estimate alone suffices except possibly for $p=2,r=1$.
In that case some positive component $b$ is repeated, so the multinomial
contains the even factor $\binom{2b}{b}$.  This supplies the missing binary
power.
\end{proof}

\begin{theorem}[A183068]\label{thm:a183068}
For every prime $p$ and all $n,r\geq1$,
\[
 a(np^r)\equiv a(np^{r-1})\pmod {p^{2r}}.
\]
\end{theorem}

\begin{proof}
Set $N=np^r$.  If $p\nmid k$, Lemma \ref{lem:carry} gives
$F(N,k)\equiv0\pmod {p^{2r}}$.  The indices divisible by $p$ are
$k=p\ell$, and Lemma \ref{lem:summand} replaces their summands by
$F(N/p,\ell)$.  Summing gives exactly $a(N/p)$ modulo $p^{2r}$.
\end{proof}

\section{Half-integral coefficient families}

For $m\geq2$, define
\[
 A_m(N)=[x^N]\left(\frac{(1+x)^m}{(1-x)^{m-2}}\right)^N.
\]
Residue extraction gives
\begin{equation}\label{eq:oddunitform}
 A_m(N)=4^N\binom{(mN-1)/2}{N}
 =\frac{2^N}{N!}\prod_{j=0}^{N-1}(mN-(2j+1)).
\end{equation}

\begin{theorem}[odd-unit block family]\label{thm:oddblock}
For $m\geq2$, $p\geq5$, and $n,r\geq1$,
\[
 A_m(np^r)\equiv A_m(np^{r-1})\pmod {p^{3r}}.
\]
\end{theorem}

\begin{proof}
Let $N=np^r$ and $M=N/p$.  In \eqref{eq:oddunitform}, the factors indexed
by odd multiples of $p$ reproduce the level $M$ product.  The remaining
quotient factors as
\[
 \frac{\binom{2N}{N}}{\binom{2M}{M}}
 \prod_{\substack{1\leq t<2N\\t\ {
m odd},\ p\nmid t}}
 \left(1-\frac{mN}{t}\right).
\]
The first factor is $1$ modulo $p^{3r}$ by \eqref{eq:LJK}.  The odd
indices form complete unit blocks modulo $p^{v_p(N)}$; their first two
reciprocal moments vanish to the precisions of Lemma \ref{lem:unit}.
Expanding the finite product proves the same cubic precision for the second
factor.  Integrality follows from the coefficient definition.
\end{proof}

The cases $m=3$ and $m=5$ prove the conjectures on \OEIS{A091527} and
\OEIS{A262732}.  Both fail at $p=3$, so the prime restriction is sharp for
these named cases.

For $a\geq3$, set
\[
 D_a(N)=\sum_{k=0}^N
 \binom{(a-1)N-k-1}{N-k}\binom{aN}{k}^2.
\]
Dixon's terminating ${}_3F_2(1)$ evaluation yields
\[
 D_a(N)=\binom{aN}{2N}
 \binom{(a+2)N/2}{N}\binom{2N}{N}
 \binom{aN/2}{N}^{-1}.
\]
Theorem \ref{thm:rbin} with $q=2$, together with \eqref{eq:LJK}, proves
\[
 D_a(np^r)\equiv D_a(np^{r-1})\pmod {p^{3r}}
 \qquad(p\geq5).
\]
The cases $a=3,5$ are \OEIS{A275652} and \OEIS{A275654}; row $a=7$ is
\OEIS{A364304}.  The finite sum supplies integrality.

\section{Symmetric-box plane partitions}

For $c,N\geq1$, let
\[
 B_c(N)=\prod_{i,j=1}^N\frac{cN+i+j-1}{i+j-1}.
\]
MacMahon's formula identifies this with the number of plane partitions in
an $N\times N\times cN$ box.

\begin{theorem}[symmetric-box Frobenius tower]\label{thm:box}
For every prime $p$ and all $c,n,r\geq1$,
\[
 B_c(np^r)\equiv B_c(np^{r-1})^p\pmod {p^{4r}}.
\]
\end{theorem}

\begin{proof}
Pair the multiplicity-$s$ and multiplicity-$(2N-s)$ factors in MacMahon's
product to obtain
\[
 B_c(N)=(c+1)^N\prod_{s=1}^{N-1}
 \left(1+\frac{c(c+2)N^2}{s(2N-s)}\right)^s.
\]
After replacing $N$ by $pN$, the factors with $p\mid s$ reproduce the
$p$-th power of the level $N$ product.  The residual unit product is
\[
 U_{p,c}(N)=\prod_{\substack{1\leq s<pN\\p\nmid s}}
 \left(1+\frac{c(c+2)p^2N^2}{s(2pN-s)}\right)^s.
\]
For odd $p$, the logarithmic linear term is a multiple of the reciprocal
sum over the units in the interval $(np^r,2np^r)$.  Pairing $t$ with
$3np^r-t$ shows that this reciprocal sum vanishes modulo $p^{2r}$.
The linear term therefore has valuation at least $4r$, and every higher
logarithmic term has at least the same valuation.

For $p=2$, complementary odd residues give the required reciprocal-block
bound.  If $c$ is even, $c(c+2)$ supplies three bonus powers.  If $c$ is
odd, the residual product is initially two powers short, but $B_c(N)$ is
even: for odd $N$, box complementation is fixed-point-free, and Legendre's
formula gives $v_2(B_c(2N))=2v_2(B_c(N))$.  Thus the square
$B_c(N)^2$ supplies the two missing powers.
\end{proof}

The cases $c=1,2,3$ prove the named conjectures on \OEIS{A008793},
\OEIS{A352656}, and \OEIS{A352657}.

\section{A spin-foam-inspired state-sum interpretation}

Spin-foam models assign local labels to the faces and edges of a
two-complex, multiply local amplitudes, and sum over compatible labelings
\cite{BaezSpinFoam,PerezSpinFoam}.  Their renormalization problem compares
amplitudes on fine and coarse discretizations \cite{Steinhaus}.  The
supercongruence towers above admit a precise arithmetic analogue of this
\emph{state-sum and coarse-graining architecture}, although they are not
models of quantum gravity.

For \OEIS{A183068}, a state at scale $N$ is a six-color occupancy vector
\[
 (k,k,k,k,N-k,N-k)
\]
with local multiplicity $F(N,k)$.  The partition function is
$a(N)=\sum_kF(N,k)$.  The coarse-graining map at the prime $p$ sends the
surviving state $k=p\ell$ to $\ell$.  States not lying in the image acquire
at least $2r$ powers of $p$ from carries, while the amplitude of a surviving
state agrees with its coarse amplitude modulo $p^{2r}$.  Thus Theorem
\ref{thm:a183068} is a local-to-global invariance statement for an
arithmetic partition function.

The product models have the same shape with no external summation.  A
microscopic label is an individual numerator or denominator factor.  The
$p$-divisible labels reproduce the coarser product exactly, and the unit
labels form complete fibers over $(\Z/p^e\Z)^\times$.  Lemma
\ref{lem:unit} says that the first two moments of every fiber vanish to the
needed $p$-adic precision.  In this language the passage
\[
 N\longmapsto pN
\]
is an arithmetic subdivision, and a congruence such as
\[
 Z(pN)\equiv Z(N)\pmod {p^{\lambda r}}
\]
or $Z(pN)\equiv Z(N)^p$ is a finite-precision renormalization fixed-point
law.  The exponent $\lambda r$ measures the depth to which fine-scale
fluctuations are invisible in the coarse observable.

This viewpoint suggests a reusable definition.  An \emph{arithmetic foam}
is a finite labeled cell system with amplitudes in $\Z_p$, together with a
surjective coarse-graining map whose nontrivial fibers have controlled
$p$-adic moments.  A state-sum supercongruence follows when
\begin{enumerate}
\item amplitudes outside the coarse image vanish to the target depth;
\item amplitudes inside the image transfer to their coarse counterparts;
      and
\item these estimates are stable under summation or multiplication.
\end{enumerate}
The carry lemma and the unit-block lemma are two concrete realizations of
the same three axioms.

The analogy stops there.  No gauge group, representation labels,
intertwiner spaces, Pachner-move invariance, semiclassical limit, or
continuum theory has been constructed.  Consequently these arithmetic
foams should be read as a physics-inspired organizational device, not as a
claim that the OEIS sequences are spin-foam quantum-gravity amplitudes.
The useful research question is narrower: whether known spin-foam and
tensor-network coarse-graining identities suggest new finite-fiber moment
cancellations, and whether those cancellations can then be proved purely
arithmetically.

\section{Compact portfolio ledger}

Table \ref{tab:ledger} records the named OEIS consequences presently
packaged as complete proofs or explicit deductions in the public
repository.  It is a review index, not a claim that every line uses the
same proof.  The four engines proved in this manuscript cover the rows
marked ``this paper.''  The remaining rows have self-contained companion
notes and exact checkers in the repository and should be audited separately
before being incorporated into a final arXiv submission.

\small
\begin{longtable}{>{\raggedright\arraybackslash}p{0.24\textwidth}
                  >{\raggedright\arraybackslash}p{0.29\textwidth}
                  >{\raggedright\arraybackslash}p{0.37\textwidth}}
\caption{Named supercongruence consequences in the review portfolio.}
\label{tab:ledger}\\
\toprule
OEIS records & Mechanism & Review status \\
\midrule
\endfirsthead
\toprule
OEIS records & Mechanism & Review status \\
\midrule
\endhead
\OEIS{A183068} & multinomial carries and scaling & Theorem \ref{thm:a183068}; all primes \\
\OEIS{A008793}, \OEIS{A352656}, \OEIS{A352657} & symmetric-box product & Theorem \ref{thm:box}; all primes \\
\OEIS{A364175}, \OEIS{A365025}--\OEIS{A365027},
\OEIS{A364515}--\OEIS{A364517} & residue-balanced rational gamma ratios & Theorem \ref{thm:gamma}; $p\geq5$ \\
\OEIS{A091527}, \OEIS{A262732} & odd unit blocks & Theorem \ref{thm:oddblock}; $p\geq5$ \\
\OEIS{A275652}, \OEIS{A275654}, \OEIS{A364304} & Dixon and half-binomial transfer & Section 5; $p\geq5$ \\
\OEIS{A061164}, \OEIS{A364506} & balanced factorial or binomial Laurent products & complete deductions from \eqref{eq:LJK} \\
\OEIS{A288470}, \OEIS{A141057} & termwise carry/Frobenius families & complete companion proofs; external review pending \\
\OEIS{A357509}, \OEIS{A357568} & binomial-quotient cancellation & complete companion proofs; external review pending \\
\OEIS{A002003}, \OEIS{A348410}, \OEIS{A351857},
\OEIS{A352373}, \OEIS{A370101}, \OEIS{A370102} & coefficient framing/Cartier transfer & complete companion proofs; external review pending \\
\OEIS{A008485}, \OEIS{A008705}, \OEIS{A270913}, \OEIS{A270919} & modular-product prime coefficients & complete companion proofs; external review pending \\
\OEIS{A255672}, \OEIS{A270922}, \OEIS{A270924} & colored Euler-product transfer & complete companion proofs; external review pending \\
\OEIS{A108625}, \OEIS{A143007}, \OEIS{A177316} & multivariate Ap\'ery transfer at the boundary prime & complete companion proofs; external review pending \\
\OEIS{A049505} & paired symmetric-plane-partition products & complete companion proof; external review pending \\
\OEIS{A357510} & odd Ap\'ery moment at primes & complete companion proof; composite extensions remain separate \\
\bottomrule
\end{longtable}
\normalsize

The repository also contains reductions and computational targets.  They
are intentionally absent from Table \ref{tab:ledger}; in particular, the
enhanced Ap\'ery defect relations, paired ordinary/shifted towers, and the
remaining global integrality problems are not declared solved.

\section{Verification and review questions}

The exact checkers use integer or rational arithmetic and are intended to
catch transcription errors, not to replace proofs.  The main commands are
\begin{verbatim}
python verification/verify_a183068.py
python verification/related/verify_rational_gamma_ratio_towers.py
python verification/related/verify_symmetric_box_plane_partitions.py
python verification/related/verify_dixon_legendre_towers.py
python verification/check_repository_integrity.py
\end{verbatim}
The rational-gamma checker performs 8,666 exact checks; the A183068,
symmetric-box, Dixon--Legendre, and repository-integrity suites are also
registered by the public repository runner.

For an efficient independent review, the highest-risk obligations are:
\begin{enumerate}
\item verify the exact cancellation identity in Theorem \ref{thm:rbin};
\item verify that residue-class balance is sufficient for the pairing in
      Theorem \ref{thm:gamma}, including negative multiplicities;
\item audit the small-prime exponent and sign in
      \eqref{eq:multiscale}, especially $p=2,r=1$;
\item audit the binary parity repair in Theorem \ref{thm:box}; and
\item check each OEIS formula and the literature-priority boundary
      independently of the proofs.
\end{enumerate}

\section*{Acknowledgments}

We thank Peter Bala for formulating the OEIS conjectures and for identifying
the superfactorial, Ap\'ery-combination, paired-tower, and
fractional-factorial directions.  We thank Paul D. Hanna for creating
\OEIS{A183068} and for encouraging review of the proposed proof.  We thank
David Radcliffe for making his integrality proof for \OEIS{A364175}
publicly available.  This research used large language models for
conjecture routing, symbolic exploration, proof criticism, and code
generation.  Every promoted claim is accompanied by a written proof and
exact reproducibility checks; machine assistance is not described as peer
review.

\begin{thebibliography}{99}

\bibitem{Bober}
J. W. Bober,
\emph{Factorial ratios, hypergeometric series, and a family of step
functions}, J. London Math. Soc. 79 (2009), 422--444;
\href{https://arxiv.org/abs/0709.1977}{arXiv:0709.1977}.

\bibitem{BaezSpinFoam}
J. C. Baez,
\emph{Spin foam models},
Class. Quantum Grav. 15 (1998), 1827--1858;
\href{https://arxiv.org/abs/gr-qc/9709052}{arXiv:gr-qc/9709052}.

\bibitem{Kazandzidis}
G. S. Kazandzidis,
\emph{Congruences on the binomial coefficients},
Bull. Soc. Math. Gr\`ece, N.S. 9 (1968), 1--12.

\bibitem{OSS}
R. Osburn, B. Sahu, and A. Straub,
\emph{Supercongruences for sporadic sequences},
Proc. Edinburgh Math. Soc. 59 (2016), 503--518;
\href{https://arxiv.org/abs/1312.2195}{arXiv:1312.2195}.

\bibitem{Radcliffe}
D. Radcliffe,
\emph{Integrality of a ratio of fractional factorials},
note linked from \OEIS{A364175}, July 2026;
\href{https://oeis.org/A364175/a364175.pdf}{OEIS PDF}.

\bibitem{PerezSpinFoam}
A. Perez,
\emph{The new spin foam models and quantum gravity},
Papers Phys. 4 (2012), 040004;
\href{https://arxiv.org/abs/1205.0911}{arXiv:1205.0911}.

\bibitem{Soundararajan}
K. Soundararajan,
\emph{Integral factorial ratios: irreducible examples with height larger
than 1}, Philos. Trans. Roy. Soc. A 378 (2020), 20180444;
\href{https://arxiv.org/abs/1906.06413}{arXiv:1906.06413}.

\bibitem{StraubLjunggren}
A. Straub,
\emph{A $q$-analog of Ljunggren's binomial congruence},
S\'em. Lothar. Combin. 64 (2011), Article B64c;
\href{https://arxiv.org/abs/1103.3258}{arXiv:1103.3258}.

\bibitem{StraubMulti}
A. Straub,
\emph{Multivariate Ap\'ery numbers and supercongruences of rational
functions}, Algebra Number Theory 8 (2014), 1985--2008;
\href{https://arxiv.org/abs/1401.0854}{arXiv:1401.0854}.

\bibitem{Steinhaus}
S. Steinhaus,
\emph{Coarse graining spin foam quantum gravity -- a review},
Front. Phys. 8 (2020), 295;
\href{https://arxiv.org/abs/2007.01315}{arXiv:2007.01315}.

\bibitem{Zudilin}
W. Zudilin,
\emph{Congruences for $q$-binomial coefficients},
Ann. Comb. 23 (2019), 1123--1135;
\href{https://arxiv.org/abs/1901.07843}{arXiv:1901.07843}.

\end{thebibliography}

\end{document}
